\textbf{(CAP-COLUNI 2021 -- Q12)} No Brasil, desde dezembro de 2018, começaram a ser usadas as placas no padrão Mercosul para os veículos 0 km, os que passaram por mudança de município, os que trocaram de categoria e aqueles com placa ilegível ou danificada.

Disponível em: <https://quatrorodas.abril.com.br/auto-servico/placa-do-mercosul-tire-suas-duvidas-e-saiba-o-que-ja-mudou-no-projeto/>. Acesso em: 20 nov. 2020.

As novas placas têm três letras seguidas; um número; uma letra; e depois mais dois números, nessa ordem, como o exemplo abaixo.

Considerando as dez primeiras letras do alfabeto e os algarismos pares, assinale a alternativa que indica CORRETAMENTE quantas placas com quatro letras distintas podem ser formadas:

\includegraphics{./figs/2021-Q12.png}

Disponível em: <http://telegrafoonline.com.br/noticias/brasil/2018/09/11/placas-padrao-mercosul-comecam-ser-utilizadas-no-brasil/>. Acesso em: 20 nov. 2020.

\begin{enumerate}
	\item 302.400
	\item 630.000
	\item 1.125.000
	\item 1.250.000
\end{enumerate}
